% Chapter 12: Use Cases, Case Studies, and Cost Scenarios
\chapter{USE CASES, CASE STUDIES, AND COST SCENARIOS}
\thispagestyle{plain}

Multi-cloud orchestration becomes more valuable when assessed in terms of practical use case applications. This chapter will present several representative examples to highlight how Nebula can be used for various needs. There are no claims about universality and perfection made, instead, it aims to illustrate the power of a unified control plane as an abstraction for decision-making, deployment choices, and financial considerations.

\section{Use Case 1: Startup SaaS Platform}


Assume a startup working on customer-facing SaaS product and having a small engineer- ing group. The goal is to develop the solution fast and clearly, minimize upfront costs, and refrain from getting hooked up with one specific cloud provider for now. 


In such a situation, Nebula can assist with experimentation by offering a standardized interface for developing VMs, object stores, regional instances, etc. Instead of having to get used to all three portals immediately, the company will be able to make standardized requests and compare results side-by-side.


The key value here is cognitive efficiency. The engineering team is typically small and consists of people who both implement the solution and provide the corresponding infrastructure. Nebula will help reduce cognitive load, facilitate onboarding, and retain knowledge of what was deployed where and why.


For this type of organization, the main benefit is cognitive efficiency. The team is usually small, and the same engineers often handle both product features and infrastructure work. Nebula reduces context switching, shortens onboarding time for new team members, and preserves a record of what was provisioned and why.

\section{Use Case 2: Enterprise Department With Microsoft Alignment}


Assume an enterprise department that heavily utilizes Microsoft-related services for their identity and collaboration needs but wants certain analytics or external-facing solutions to be deployed in other clouds. For example, there might be an interest in leveraging AWS/GCP as opposed to Azure because the workload is better suited to particular providers. In such a scenario, Nebula can serve as a tool to keep Azure-based policy framework while still being able to manage other services in controlled manner.


In such an environment, multi-cloud orchestration is mostly about the policy-boundary nature as opposed to merely being convenient. The importance of onboarding, tag enfor- cing, restriction to regions, and auditability cannot be underestimated. However, having the same dashboard remains important because the department should know where resources are hosted regardless of the need for everyone to be proficient in each provider's console.
e.

\section{Use Case 3: Research Laboratory or Academic Institution}


Another scenario in which Nebula will work efficiently is when there are periodic demands for a computing cluster, storage facility, and reproducibility of the infrastructures for academic assignments or experiments. In such cases, budget visibility and repeatability will help since new people join the group and leave after some time, thereby making infrastructure information fragment among different groups of students.


The ability of Nebula to document, observe, and reuse workflows suits such scenarios well. This is because the control plane allows a technical lead to define a series of standardized procedures for setup and provide infrastructure templates for them to follow. Moreover, the technical leads will be able to see all the components and what has been set up using various providers.

\section{Use Case 4: Seasonal E-Commerce Expansion}


In the case where e-commerce sites experience seasonal demands, it might be desirable to maintain one cloud provider while being ready to use a second one in case the first one fails. In most times, not all resources will be used at any time depending on the level of activity. However, during high activity levels, it might be desirable to use the second cloud as well.


In such scenarios, the team must know how to distinguish between primary and backup resources, how resource inventory can be kept synchronized, and how faults can be analyzed quickly. This is one of the challenges the control plane can address by allowing task tracking and centralized inventory viewing.

\section{Illustrative Cost Scenario Comparison}
The following table provides an example of the kind of workload-based decision-making the control plane will allow. The numbers used in the table below are illustrative and not vendor quotes.

\begin{table}[H]
\centering
\caption{Illustrative Cost and Fit Comparison by Scenario}
\label{tab:cost_scenario}
\begin{tabular}{|p{4.0cm}|p{2.0cm}|p{2.0cm}|p{3.0cm}|}
\hline
\textbf{Scenario} & \textbf{Primary Goal} & \textbf{Suggested Cloud Bias} & \textbf{Nebula Value} \\ \hline
Student development sandbox & low cost & burstable mixed cloud & one workflow for repeat setup \\ \hline
Enterprise internal app & identity alignment & Azure-first & governance and inventory unification \\ \hline
Data-heavy analytics prototype & fast experimentation & GCP-first & compare supporting infrastructure centrally \\ \hline
Broad service integration platform & service breadth & AWS-first & consistent provisioning and failure handling \\ \hline
Disaster recovery standby & resilience & dual-provider & coordinated fallback orchestration \\ \hline
\end{tabular}
\end{table}

\section{Case Study: Failed Provisioning With Guided Recovery}
One of the most common examples of how useful Nebula is in practical terms would be an error at the moment of a time-bound release, such as a deployment failing due to a permissions error in AWS, which is detected via Terraform. In a traditional case, the operator has to decipher the log output of the provider, determine what permissions are missing, access the IAM console separately and find out if the issue is related to credentials, scope, or regions.


In the case of Nebula, on the contrary, the operator would be presented with the record of the task completed with a status indicating a failure, the log output would be saved in one place, while the AI copilot will analyze the log and make suggestions as to why the deployment failed. Regardless of the accuracy of these recommendations, in most cases, narrowing down the search space will save additional time for the operator, which is crucial in stressful situations.


\section{Case Study: Inventory Drift Detection}


The other issue, which is common in any cloud provider's environment, is called inventory drift. It usually happens when resources are manually deleted via a provider portal or modified without following the automated deployment path. That leads to a situation when local inventory and real resources in the cloud have different state. With Nebula's periodic synchronization procedure, it is easier to bring the two back into agreement by reconciling provider state with the database.


In case a company works with several clouds simultaneously, such functionality will pay off many times over by providing accurate information about cloud estate in one central place.


\section{Organizational Benefits Beyond Provisioning}
The long-term value of a system like Nebula is not limited to provisioning speed. It also creates benefits in:
\begin{itemize}
    \item \textbf{training}: new engineers learn one operational workflow first,
    \item \textbf{documentation}: a single report or manual can describe cross-cloud actions consistently,
    \item \textbf{governance}: audit trails and approval patterns become easier to define,
    \item \textbf{support}: failures are easier to triage through centralized task records,
    \item \textbf{planning}: teams can compare provider choices from a shared interface.
\end{itemize}

\section{Limitations of Scenario-Based Evaluation}

The above examples are meant to be representative rather than exhaustive. Actual organizations might face other complications, such as different security boundaries, network topology issues, data privacy concerns, or procurement regulations that our academic prototype cannot account for. Nonetheless, we can draw several insights from analyzing these scenarios and relate them back to the design considerations in the earlier chapters

\section{Discussion}

The above scenarios demonstrate the primary claim made throughout this report. Multi-cloud operations are as much about orchestration and mental overhead as they are about connecting to cloud APIs. With its ability to consolidate intent and reduce operational fragmentation, Nebula helps achieve better productivity, better governance, and better reproducibility. In particular, it allows startups to get their feet wet more quickly, helps enterprises establish governance standards, and enables researchers and academics to conduct their work in a more systematic fashion.

\newpage
